% =====================================================================
%  Project description: UDA for precipitation super-resolution
%  Template: arxiv preprint style
%  https://www.overleaf.com/latex/templates/style-and-template-for-
%  preprints-arxiv-bio-arxiv/pkzcrhzcdxmc
% =====================================================================

\documentclass{article}

\usepackage{arxiv}

\usepackage[utf8]{inputenc}
\usepackage[T1]{fontenc}
\usepackage{hyperref}
\usepackage{url}
\usepackage{booktabs}
\usepackage{amsfonts}
\usepackage{amsmath}
\usepackage{amssymb}
\usepackage{amsthm}
\usepackage{nicefrac}
\usepackage{microtype}
\usepackage{graphicx}
\usepackage{natbib}

\newtheorem{assumption}{Assumption}
\newtheorem{proposition}{Proposition}

\title{Diagnosing Failure Modes of Unsupervised Domain Adaptation in Precipitation Super-Resolution}

\author{
    Filippo Quarenghi \\
    Faculty of Geosciences and Environment (FGSE)\\
    Experise Center for Climate Extremes (ECCE)\\
    University of Lausanne \\
    Lausanne, Switzerland \\
    \texttt{filippo.quarenghi@unil.ch} \\
        \And
    Tom Beucler \\
    Faculty of Geosciences and Environment (FGSE)\\
    Experise Center for Climate Extremes (ECCE)\\
    University of Lausanne\\
    Lausanne, Switzerland\\
}

\date{}

\begin{document}
\maketitle

\begin{abstract}
Deep-learning models for precipitation super-resolution lose substantial
skill when deployed in geographic regions different from those used for
training, a problem the recent RainShift benchmark
formalises~\citep{harder2025rainshift}. We present a diagnostic framework
for evaluating unsupervised domain adaptation (UDA) in this setting,
grounded in two complementary theoretical bounds for the target risk:
the optimal-transport bound of \citet{redko2017theoretical} and the
$\mathcal{H}\Delta\mathcal{H}$-distance bound of \citet{bendavid2010theory}.
The framework decomposes the transfer gap into an addressable
covariate-shift component and an irreducible conditional-shift component,
estimates both pre-training, via multi-source consistency on supported
inputs and joint-marginal alignment residuals, and post-hoc, via oracle
joint training, and proposes a novel regime-decomposed irreducible risk
$\lambda_k$ that localises conditional shift to specific physical
regimes of input space. Across the full 3$\times$3 matrix the oracle recovers
the target's own in-domain risk to within $0.02$ standardised MSE on
every pair, so the class-restricted irreducible risk is negligible: one
network fits both conditionals at once. The transfer gap is therefore
reducible in principle yet not identifiable from unlabelled target
data, and two direct diagnostics we introduce localise the obstruction
to a difference in $P(y\mid x)$ rather than to model capacity. No
tested covariate-alignment method closes it. We propose a 3$\times$3 transfer matrix experiment across three
RainShift regions, applied to a deterministic UNet and a probabilistic
Adaptive Flow Matching model, with eight UDA methods organised by where
they act in the pipeline. The methodology generalises beyond climate
downscaling to any transfer-learning setting where the relative
dominance of covariate versus conditional shift must be diagnosed before
committing compute to adaptation.
\end{abstract}

\keywords{unsupervised domain adaptation \and precipitation
super-resolution \and optimal transport \and climate downscaling \and
distribution shift diagnostics \and irreducible risk}

% =====================================================================
\section{Introduction}
% =====================================================================

High-resolution precipitation observations are concentrated in the
Global North, while many of the regions where high-resolution
precipitation forecasts have the largest socioeconomic value lie in
data-sparse parts of the Global South. Deep-learning super-resolution
models trained on rich-data regions and deployed in sparse-data regions
are the natural route to bridging this gap, but the recent RainShift
benchmark shows that zero-shot geographic transfer degrades skill by up
to 30\% depending on architecture and region pair~\citep{harder2025rainshift}.
The RainShift paper evaluates several in-distribution baselines together
with quantile mapping on total precipitation as a simple alignment
technique; it does not benchmark unsupervised domain adaptation
systematically, nor diagnose which aspects of the transfer gap any
adaptation method could possibly close.

This project addresses the diagnostic gap. We construct an evaluation
framework that, given a source--target domain pair and a labelled source
sample, predicts before any UDA training whether covariate-shift
adaptation can help, quantifies the irreducible conditional-shift
component via oracle joint training, identifies in which physical
regimes of input space the conditional shift lives, and stratifies
post-hoc method comparison accordingly. We then benchmark a
representative set of UDA methods on a $3\times 3$ matrix of RainShift
transfers using two architectures: a deterministic dual-encoder UNet
and a probabilistic Adaptive Flow Matching
model~\citep{fotiadis2025adaptive}.

The framework is built on three observations. First, the theoretical
literature on UDA decomposes target risk into a source-risk term, a
divergence between distributions, and an irreducible joint-risk term;
the third encodes conditional shift and cannot be addressed by methods
that operate on marginals
alone~\citep{bendavid2010theory,redko2017theoretical}. Second,
identifiability of conditional shift from unlabelled target data is
impossible without additional assumptions; different UDA method classes
operationalise different assumptions, and the choice of method should
be guided by which assumption is consistent with the data. Third, the
irreducible component has structure that the scalar quantity $\lambda$
does not capture: it can be decomposed across physically meaningful
regions of input space, providing actionable guidance about where the
transfer is irreducibly hard.

% =====================================================================
\section{Theoretical foundations}\label{sec:theory}
% =====================================================================

\subsection{Setup and notation}

Let $\mathcal{X}\subset\mathbb{R}^{C\times H\times W}$ be the input
space of $C=9$ atmospheric channels at coarse resolution, including
convective available potential energy, total column water vapour,
surface pressure, two wind components, and the inputs documented in
RainShift. Let $\mathcal{Y}\subset\mathbb{R}^{1\times H'\times W'}$ be
the high-resolution precipitation field. A domain is a joint
distribution $\mathcal{D}$ over $\mathcal{X}\times\mathcal{Y}$ with
input marginal $\mu$ and conditional $P(y\mid x)$.

We are given a source domain $\mathcal{D}_S$ and a target domain
$\mathcal{D}_T$. Training data consist of a labelled source sample
$\mathcal{S}=\{(x_i^S,y_i^S)\}_{i=1}^{n_S}$ and an unlabelled target
sample $\mathcal{T}=\{x_j^T\}_{j=1}^{n_T}$. For a hypothesis class
$\mathcal{H}$ and the squared loss in standardised log-precipitation
space, the target risk
$\mathcal{R}_T(f)=\mathbb{E}_{(x,y)\sim\mathcal{D}_T}\!\left[\|f(x)-y\|^2\right]$
is the quantity to minimise under the constraint that target labels
are not observed during training. The source risk $\mathcal{R}_S(f)$
is defined analogously.

\subsection{The optimal-transport bound and its absorbed constants}
\label{sec:redko}

For a deterministic regressor and a Lipschitz loss-composition,
\citet{redko2017theoretical} establish
\begin{equation}
\mathcal{R}_T(f) \le \mathcal{R}_S(f) + K \cdot W_1(\hat\mu_S, \hat\mu_T)
+ \lambda + \mathrm{Conf}(n_S, n_T, \delta),
\label{eq:redko_full}
\end{equation}
holding with probability $1-\delta$, where four quantities deserve
explicit comment because they are routinely absorbed in applied work:

\paragraph{The Lipschitz constant $K$.} $K$ is the Lipschitz constant
of $x \mapsto \ell(f(x), y(x))$ over the support of $\mu_S \cup \mu_T$.
For deep neural networks, $K$ admits the spectral upper bound
$K \le L_\ell \cdot \prod_\ell \|W_\ell\|_2$, where $L_\ell$ is the
Lipschitz constant of the loss with respect to its first argument and
$\|W_\ell\|_2$ are layer-wise spectral norms~\citep{bartlett2017spectrally}.
For realistically-trained models this product is typically $10^2$ to
$10^4$ and is dominated by a handful of layers. The bound
\eqref{eq:redko_full} is therefore not numerically tight; we absorb $K$
into the analysis and treat the bound as a decomposition device rather
than a numerical inequality.

\paragraph{The finite-sample Wasserstein gap.} The empirical
$W_1(\hat\mu_S, \hat\mu_T)$ converges to the population $W_1$ at rate
$\mathcal{O}(n^{-1/d})$ in input dimension $d$~\citep{fournier2015rate}.
For a 9-dimensional input space at sample sizes of order $10^3$, the
estimator carries non-negligible bias; we mitigate by computing $W_1$
per channel on standardised distributions and aggregating, accepting
that this loses inter-channel coupling but improves the convergence
rate to $\mathcal{O}(n^{-1/2})$ per channel.

\paragraph{The confidence term.} $\mathrm{Conf}(n_S, n_T, \delta)$ scales
as $\mathcal{O}(\sqrt{d_\mathcal{H}/n_S + d_\mathcal{H}/n_T})$ where
$d_\mathcal{H}$ is a hypothesis-class complexity measure. For
deep models this term is worst-case vacuous; tighter bounds via
PAC-Bayes \citep{germain2013pac,dziugaite2017computing} replace
$d_\mathcal{H}$ with an effective dimension that depends on the
realised solution and can be empirically non-vacuous. We do not pursue
PAC-Bayes bounds in this work; we note them as future direction
(Section~\ref{sec:future}).

\paragraph{The irreducible joint risk $\lambda$.} Defined as
\begin{equation}
\lambda = \min_{h \in \mathcal{H}}
\big[\mathcal{R}_S(h) + \mathcal{R}_T(h)\big],
\label{eq:lambda_def}
\end{equation}
the irreducible joint risk depends on $\mathcal{H}$, not on the data
alone. The class-restricted definition is what oracle joint training
estimates and is the operationally meaningful quantity. The
Bayes-optimal version (over all measurable functions) is strictly
smaller but not empirically estimable; we are explicit throughout that
all $\lambda$-related statements are class-restricted.

For ranking and decomposition purposes we use the simplified form
\begin{equation}
\mathcal{R}_T(f) \lesssim \mathcal{R}_S(f) + W_1(\hat\mu_S, \hat\mu_T) + \lambda,
\label{eq:redko_simplified}
\end{equation}
in which $K$, the finite-sample correction, and the confidence term
are absorbed into the ``$\lesssim$'' notation. This is the form that
motivates the diagnostic decomposition into addressable and irreducible
components.

\subsection{The $\mathcal{H}\Delta\mathcal{H}$-distance bound and its
empirical estimator}
\label{sec:bendavid}

An alternative line of theoretical analysis~\citep{bendavid2010theory}
bounds the target risk in terms of the $\mathcal{H}\Delta\mathcal{H}$-%
distance:
\begin{equation}
\mathcal{R}_T(f) \le \mathcal{R}_S(f)
+ \tfrac{1}{2} d_{\mathcal{H}\Delta\mathcal{H}}(\mu_S, \mu_T) + \lambda,
\label{eq:bendavid}
\end{equation}
where
\begin{equation}
d_{\mathcal{H}\Delta\mathcal{H}}(\mu_S, \mu_T)
= 2 \sup_{h, h' \in \mathcal{H}} \Big|
\mathbb{P}_{x\sim\mu_S}[h(x) \neq h'(x)]
- \mathbb{P}_{x\sim\mu_T}[h(x) \neq h'(x)]
\Big|
\end{equation}
quantifies the maximum disagreement between two hypotheses in
$\mathcal{H}$ across the two domains. Two features distinguish this
bound from the optimal-transport bound \eqref{eq:redko_simplified}.

First, the shift quantity is defined \emph{relative to the hypothesis
class} used for adaptation, not relative to arbitrary measurable
functions. This is tighter and more honest about what the model can
actually exploit.

Second, $d_{\mathcal{H}\Delta\mathcal{H}}$ admits a natural empirical
estimator: the accuracy of a domain discriminator trained to separate
source and target features. This is the theoretical anchor for
adversarial domain adaptation~\citep{ganin2016domain}.

The framework exploits a related observation. When multiple
source-trained models $\{h_1, \ldots, h_K\}$ are evaluated on a target
input, their disagreement provides a sampled estimator of
$d_{\mathcal{H}\Delta\mathcal{H}}$ where the two hypotheses are drawn
from different source training distributions rather than from random
initialisations~\citep{mansour2008domain,lakshminarayanan2017simple}.
High inter-source disagreement on supported target inputs implies that
the multi-source ensemble cannot be made domain-invariant by any
single $\phi$ in the class, which provides theoretical justification
for the multi-source consistency diagnostic introduced in
Section~\ref{sec:framework}.

\subsection{Identifiability and the assumptions that close it}
\label{sec:identifiability}

Pure conditional shift is not identifiable from unlabelled target data
alone: if $\mu_T = \mu_S$ exactly, no statistic of $\{x_j^T\}$ can
distinguish $P_T(y\mid x)$ from $P_S(y\mid x)$. Practical UDA methods
close this identifiability gap through additional assumptions; we
enumerate four families and map each UDA method class to the
assumption it operationalises.

\begin{assumption}[Covariate-shift]
\label{ass:cov}
$P_S(y\mid x) = P_T(y\mid x)$, with shift confined to input marginals.
Under this assumption, importance-weighted empirical risk minimisation
is consistent~\citep{sugiyama2007direct}. Methods that operationalise
it: input-level alignment (QM, FDA, joint Sinkhorn-OT), which align
$\mu_S$ to $\mu_T$ on the assumption that the conditional is invariant.
\end{assumption}

\begin{assumption}[Target-shift]
\label{ass:tgt}
$P_S(x\mid y) = P_T(x\mid y)$, with shift in the label marginal
$P(y)$~\citep{lipton2018detecting}. For continuous regression with
structured spatial outputs (precipitation fields), this assumption
does not map cleanly to a single shift quantity and is not
operationalised by any method in our taxonomy.
\end{assumption}

\begin{assumption}[Shared invariant representation]
\label{ass:inv}
There exists a feature map $\phi$ such that $P_S(y\mid \phi(x))
= P_T(y\mid \phi(x))$. This is an assumption on the hypothesis class
admitting a domain-invariant representation. Methods that operationalise
it: feature-level alignment (CORAL, MMD, multi-scale MMD, DANN), which
search for $\phi$ within the model that satisfies the assumption.
\end{assumption}

\begin{assumption}[Causal sufficiency]
\label{ass:causal}
The observed inputs $x$ form a complete set of causes of $y$, so the
conditional $P(y\mid x)$ is mechanistically stable across
domains~\citep{scholkopf2012causal}. For precipitation, this holds at
the level of fundamental physics (Clausius-Clapeyron, convective
thermodynamics) but is violated by unresolved sub-grid processes that
act as omitted causal variables. The irreducible risk $\lambda$ in
\eqref{eq:lambda_def} is, in this sense, a measure of the causal
incompleteness of the input set. Methods that operationalise this
assumption, combined with physical priors: physical-consistency
losses, spectral PSD matching.
\end{assumption}

The framework's pre-training diagnostics in
Section~\ref{sec:framework} are designed to test which of
Assumptions~\ref{ass:cov}, \ref{ass:inv}, and \ref{ass:causal} is
consistent with the data, rather than committing to one a priori:
the joint-marginal alignment residual tests whether covariate-shift
adaptation has headroom (Assumption~\ref{ass:cov}); multi-source
consistency tests whether a shared invariant representation is
plausible (Assumption~\ref{ass:inv}); the physical-prior check tests
causal completeness (Assumption~\ref{ass:causal}).

\subsection{Regime-decomposed irreducible risk}
\label{sec:lambda_k}

A scalar $\lambda$ tells us whether transfer is irreducibly hard, but
not \emph{where} the hardness lives. For precipitation downscaling, the
practitioner-relevant question is which physical regimes of input space
contribute most to $\lambda$ and which are addressable. We propose a
regime-decomposed irreducible risk.

Let $\{\mathcal{X}_1, \ldots, \mathcal{X}_K\}$ be a partition of input
space into physical regimes computable from inputs alone (for example,
by joint convective-available-potential-energy and column-water bins,
or by Köppen climate classes). Define the regime-conditional
irreducible risk
\begin{equation}
\lambda_k = \min_{h \in \mathcal{H}}
\Big[
\mathbb{E}_S\!\left[\ell(h(X), Y) \mid X \in \mathcal{X}_k\right]
+ \mathbb{E}_T\!\left[\ell(h(X), Y) \mid X \in \mathcal{X}_k\right]
\Big].
\label{eq:lambda_k}
\end{equation}

\begin{proposition}[Regime decomposition]
\label{prop:decomp}
Under the assumption that the optimal hypothesis is regime-stable
(the inner minimiser of \eqref{eq:lambda_k} is independent of $k$),
\begin{equation}
\lambda = \sum_{k=1}^K \mu_T(\mathcal{X}_k)\, \lambda_k.
\end{equation}
Without the regime-stability assumption, the right-hand side is an
upper bound:
\begin{equation}
\lambda \le \sum_{k=1}^K \mu_T(\mathcal{X}_k)\, \lambda_k.
\end{equation}
\end{proposition}

The regime decomposition gives the framework two properties the scalar
$\lambda$ does not. First, it identifies \emph{where} in input space
the conditional shift dominates, so practitioners know which physical
regimes to flag for additional data acquisition or model refinement.
Second, it provides a falsifiable prediction: if regime $k$ has
$\lambda_k$ large, no covariate-shift UDA method should improve target
performance specifically on regime $k$. The physical-regime
stratification of UDA results in Section~\ref{sec:physical} tests this
prediction directly.

We estimate $\lambda_k$ empirically by per-regime oracle joint
training: train the joint model on the union of source and target
labels, then evaluate on the source and target test sets restricted to
$\mathcal{X}_k$. This is the natural per-regime version of the scalar
oracle estimate. To our knowledge, the regime-decomposed irreducible
risk has not been studied for regression with physically meaningful
continuous partitions. The closest published direction is the
subpopulation-shift line in benchmark
evaluation~\citep{koh2021wilds}, which considers discrete subgroup
shift for classification but does not address regression or continuous
physical regimes.

\subsection{Mapping bounds to UDA method classes}
\label{sec:mapping}

The two bounds (Sections~\ref{sec:redko}, \ref{sec:bendavid}) and the
four identifiability assumptions (Section~\ref{sec:identifiability})
induce a mapping from method class to the operative theoretical
quantity, summarised in Table~\ref{tab:mapping}. The framework's
diagnostic-to-method pipeline is: pre-training diagnostics identify
which assumption is consistent with the data, the regime-decomposed
$\lambda_k$ identifies which input regimes admit adaptation, and the
method class consistent with the surviving assumption is selected for
trial.

\begin{table}[h]
\centering
\begin{tabular}{llll}
\toprule
Method class & Assumption & Operative bound & Diagnostic signal \\
\midrule
Input alignment & Ass.~\ref{ass:cov} & Redko (marginal $W_1$)
& Low joint-marginal alignment residual \\
Feature alignment & Ass.~\ref{ass:inv} & Ben-David ($\mathcal{H}\Delta\mathcal{H}$)
& Low multi-source disagreement \\
Output / physical & Ass.~\ref{ass:causal} & Physical consistency
& Physical-prior check passes \\
Joint distribution & None alone & Joint $W_1$
& All single-assumption diagnostics fail \\
\bottomrule
\end{tabular}
\caption{Mapping from UDA method class to the identifiability assumption
it operationalises, the theoretical bound that applies, and the
pre-training diagnostic signal that supports its use.}
\label{tab:mapping}
\end{table}

When all single-assumption diagnostics signal failure, only
joint-distribution methods (JDOT)~\citep{courty2017jdot} or supervised
target adaptation (few-shot) can close the gap. This is the scenario
the europe\_west $\to$ melanesia case study currently appears to be in.

% =====================================================================
\section{Methodology}
% =====================================================================

\subsection{Experimental setup}

The benchmark consists of a $3\times 3$ transfer matrix over three
RainShift regions chosen to span a difficulty axis: europe\_west
(mid-latitude maritime), horn\_of\_africa (tropical semi-arid), and
melanesia (tropical maritime convective). The matrix yields nine
evaluation cells comprising three self-evaluations on the diagonal and
six off-diagonal transfer pairs. The diagonal cells establish in-domain
baselines; off-diagonal cells allow testing of (a) asymmetry (S$\to$T
versus T$\to$S), (b) lateral transfer (tropical to tropical), and (c)
the dependence of UDA effectiveness on a measurable shift severity.

For each of the six transfer pairs we run every UDA method with three
random seeds. With two architectures and eight methods, the experiment
totals approximately 250 training runs.

\subsection{Architectures}

The deterministic baseline is a dual-encoder UNet with nine
atmospheric input channels in the dynamic branch and two static
channels (topography and land-sea mask) in the static branch. Training
loss is mean squared error in standardised log-precipitation space.

The probabilistic baseline is the Adaptive Flow Matching (AFM) model
of~\citet{fotiadis2025adaptive}, which couples an encoder producing a
deterministic prediction $\mu$ with a flow network parameterising
optimal-transport probability paths~\citep{lipman2023flow} between a
perturbed prior and the target distribution. The encoder architecture
is identical to the deterministic UNet, so feature-level UDA losses
apply identically across the two models, providing a clean controlled
comparison between deterministic and probabilistic super-resolution
under shift. The flow network's adaptive noise scaling $\sigma_z$,
updated via an exponential moving average of encoder prediction error,
provides a partial mechanism for conditional-variance adaptation; we
test the calibration of this mechanism through spread-skill ratios on
the target.

A two-stage CorrDiff-style decomposition~\citep{mardani2024corrdiff}
was considered and rejected for this paper: applying UDA losses only
to the deterministic stage cannot guarantee that the diffusion
corrector, trained on source-domain residuals, captures the correct
target-domain conditional variance. AFM's end-to-end training and
$\sigma_z$ adaptation provide a more principled handle on this
question.

\subsection{UDA methods, organised by where they act}

We benchmark eight methods grouped by their point of intervention in
the model pipeline, with each method mapped to its operative
assumption (Table~\ref{tab:mapping}):

\paragraph{Input-level alignment (Ass.~\ref{ass:cov}).} Quantile
mapping per channel as in the RainShift baseline, and a joint Sinkhorn
optimal-transport alignment~\citep{cuturi2013sinkhorn} that respects
inter-channel correlations. We additionally include Fourier domain
adaptation~\citep{yang2020fda} as a negative-control baseline, since
preliminary results indicate that swapping low-frequency Fourier
amplitudes between atmospheric domains corrupts climatological mean
state rather than implementing a style/content separation.

\paragraph{Feature-level alignment (Ass.~\ref{ass:inv}).} Deep
CORAL~\citep{sun2016deep} matches second-order feature statistics;
multi-scale maximum mean discrepancy~\citep{gretton2012kernel,long2015learning}
matches kernel embeddings of feature distributions at multiple encoder
levels with median-heuristic bandwidths; adversarial alignment via
DANN~\citep{ganin2016domain} uses a gradient-reversal-layer domain
discriminator and is the most direct empirical estimator of the
$\mathcal{H}\Delta\mathcal{H}$ shift. For AFM, all feature-level
methods are applied to the encoder's intermediate features only.

\paragraph{Joint-distribution alignment.} Joint distribution optimal
transport~\citep{courty2017jdot} aligns both input marginals and
predicted conditionals simultaneously, dropping the
single-assumption restriction by working with the joint distribution
directly.

\paragraph{Test-time alignment.} Adaptive batch
normalisation~\citep{li2018adabn} replaces batch-normalisation running
statistics with target-domain statistics at inference.

\paragraph{Output-level alignment (Ass.~\ref{ass:causal}).} A radially
averaged power spectral density matching loss penalises differences
between source and target prediction spectra, providing a physically
motivated alignment of the spatial correlation structure of predicted
precipitation.

% =====================================================================
\section{Diagnostic framework}\label{sec:framework}
% =====================================================================

The framework is composed of three layers that together estimate where
in the bound \eqref{eq:redko_simplified} the transfer gap lives, both
before and after UDA training, and which identifiability assumption
from Section~\ref{sec:identifiability} is consistent with the data.

\subsection{Pre-training diagnostics}

Three surrogate diagnostics, each tied to one identifiability
assumption, predict the addressable budget without training any UDA
method:

\textit{Joint-marginal alignment residual} (tests
Assumption~\ref{ass:cov}). The joint Sinkhorn-OT alignment of source to
target inputs is fitted and the residual $W_1$ measured. If the
residual is non-negligible, Assumption~\ref{ass:cov} is violated and
covariate-shift adaptation has an upper-bounded ceiling. If the
residual is near zero but the transfer gap persists, the residual gap
must be attributed to conditional shift.

\textit{Multi-source consistency on supported inputs} (tests
Assumption~\ref{ass:inv} and provides an empirical estimator of
$d_{\mathcal{H}\Delta\mathcal{H}}$). Multiple source-trained models, one
per RainShift training region, are evaluated on each target input.
Disagreement among models conditional on high coverage of the input
across sources is interpreted as evidence against the existence of a
shared invariant representation in $\mathcal{H}$.

\textit{Physical-prior consistency check} (tests
Assumption~\ref{ass:causal}). Two-variable slices of input space (e.g.
column water vapour against convective available potential energy) are
compared across source and target. Disagreement in physically
meaningful regions of input space, where known precipitation mechanisms
differ, indicates that the observed inputs do not form a sufficient
causal set, so the conditional $P(y\mid x)$ is unstable.

\subsection{Oracle-grounded budget estimation}

For each transfer pair, a joint-training model is trained on the union
of source and target labels and evaluated on each domain. The
resulting $\widehat{\lambda}=\mathcal{R}_S(h_{\text{joint}}) +
\mathcal{R}_T(h_{\text{joint}})$ is a class-restricted upper bound on
$\lambda$, and
$\mathcal{R}_T(f_S)-\mathcal{R}_T(h_{\text{joint}})$ is the addressable
budget. The fraction of this budget recovered by each unsupervised UDA
method is the cleanest single quantity to report per method.

For the regime-decomposed analysis we additionally compute the
per-regime $\widehat{\lambda}_k$ by evaluating the joint model on the
test sets stratified by physical regime, producing a vector of
irreducible risks across input space.

\subsection{Post-hoc analysis with explicit gap decomposition}

For each (source, target, method) triple we report source MSE, target
MSE, and the source-target gap separately. We use the normalised
generalisation gap $\mathrm{NGG}(S,T)=(\mathcal{R}_T-\mathcal{R}_S)
/(W_1+\epsilon)$ for within-architecture rankings, but report it
alongside source-MSE inflation to expose methods that achieve apparent
target improvement through regression-to-mean rather than genuine
adaptation. Preliminary results show this distinction is necessary in
practice: feature-moment-matching methods can reduce the
target--source gap without reducing target error in absolute terms.

% =====================================================================
\section{Transfer gaps and the addressable budget}\label{sec:budget}
% =====================================================================

We report the full $3\times3$ matrix for the deterministic UNet with
source-only and oracle joint training. Two conventions matter for
reading the numbers. First, standardised MSE is computed in the
\emph{source} domain's normalised space, because that is the space the
model is trained in; consequently \texttt{mse\_std} is comparable
\emph{within} a transfer pair but not across pairs with different
sources. Second, to compare across pairs we convert each target risk
into the target domain's own normalisation using the ratio of
log-precipitation standard deviations, which is recoverable because
each domain appears as a target under two different source
normalisations. The two independent recoveries agree to within 2\%
($\sigma_{\text{horn}}/\sigma_{\text{eu}} = 0.729$ vs $0.723$;
$\sigma_{\text{mel}}/\sigma_{\text{eu}} = 1.256$ vs $1.279$), which is
itself a consistency check on the evaluation pipeline.

\begin{table}[t]
\centering
\small
\begin{tabular}{lccccc}
\toprule
pair & $\mathcal{R}_T(f_S)$ & $\mathcal{R}_T(h_{\text{joint}})$ &
in-domain & transfer gap & $\widehat{\lambda}$ \\
\midrule
horn $\to$ europe\_west      & $0.958$ & $0.529$ & $0.523$ & $0.435$ & $0.005$ \\
melanesia $\to$ europe\_west & $0.872$ & $0.542$ & $0.523$ & $0.349$ & $0.019$ \\
europe\_west $\to$ horn      & $0.815$ & $0.635$ & $0.631$ & $0.184$ & $0.004$ \\
europe\_west $\to$ melanesia & $0.770$ & $0.621$ & $0.618$ & $0.152$ & $0.002$ \\
horn $\to$ melanesia         & $0.723$ & $0.641$ & $0.618$ & $0.105$ & $0.023$ \\
melanesia $\to$ horn         & $0.717$ & $0.633$ & $0.631$ & $0.086$ & $0.002$ \\
\bottomrule
\end{tabular}
\caption{Transfer matrix in standardised MSE, converted to the
\emph{target} domain's normalisation so that rows are comparable.
$\mathcal{R}_T(f_S)$ is the source-only model on the target;
$\mathcal{R}_T(h_{\text{joint}})$ is the oracle trained on both label
sets; in-domain is the target's own source-only model. The transfer gap
is $\mathcal{R}_T(f_S)$ minus in-domain, and
$\widehat{\lambda} = \mathcal{R}_T(h_{\text{joint}})$ minus in-domain is
the class-restricted irreducible excess.}
\label{tab:budget}
\end{table}

\paragraph{The irreducible risk is negligible.}
On every pair the oracle recovers the target's own in-domain risk to
within $0.002$--$0.023$ standardised MSE, i.e.\ $\widehat{\lambda}
\approx 0$. A single UNet trained on both label sets fits both
conditionals essentially as well as two separately trained models. This
is the most consequential number in the study and it rules out the
trivial explanation for adaptation failure: the hypothesis class is not
the binding constraint, and there is no capacity-driven irreducible
term to blame.

This sharpens the claim the framework can make. The transfer gap is
\emph{reducible in principle} --- $\widehat{\lambda}\approx 0$ exhibits
a hypothesis in $\mathcal{H}$ that solves both domains simultaneously
--- but it is not \emph{identifiable} from unlabelled target data.
Labels locate that hypothesis; marginal alignment does not. The failure
is one of identifiability in the sense of
Section~\ref{sec:identifiability}, not of expressivity. We regard this
as a stronger and more precise statement than an appeal to a large
irreducible $\lambda$, and it is the statement the remaining
experiments test.

\paragraph{Transfer is strongly asymmetric.}
Direction matters more than pair identity. Transfer \emph{into}
europe\_west costs $0.435$ and $0.349$, while transfer into either
tropical domain costs $0.086$--$0.184$: a four- to fivefold difference.
The asymmetry is not visible in mean absolute error, which is nearly
insensitive on the dry midlatitude target (in-domain $0.088$~mm versus
$0.095$~mm under transfer, against a standardised-MSE gap of $0.435$),
so we report standardised MSE as the headline metric and treat MAE as a
secondary summary. All models additionally carry a negative
precipitation bias ($-0.05$ to $-0.29$~mm), the expected
under-dispersion signature of an MSE objective in log space.

\paragraph{The diagnostics predict the ordering.}
The matched-input diagnostic of Section~\ref{sec:conditional} ranks the
pairs by conditional discrepancy without training any adaptation
method, and that ranking matches the measured transfer difficulty.
horn-of-africa $\leftrightarrow$ melanesia carries the smallest excess
over the within-domain null ($14.5\times$) and the two \emph{smallest}
transfer gaps ($0.086$, $0.105$); europe\_west $\leftrightarrow$ horn
carries the largest excess ($86.9\times$) and gaps of $0.184$ and
$0.435$ depending on direction. This is the claim the framework is
built to support: a quantity computable from unlabelled inputs and
observed targets alone ranks transfer difficulty correctly.

We state the corollary plainly: on this domain triple no pair offers
covariate alignment both a theoretical licence and substantial
headroom. The pair with the mildest conditional shift also has the
smallest gap to close, and the pairs with the largest gaps are those
where the conditional differs most or where the input supports barely
overlap. This is a property of this domain triple rather than a general
result, and it is why the benchmark's informative comparison is between
the joint-distribution method and the marginal-alignment family, not
among the latter.

\paragraph{Revised budget estimate.}
For europe\_west $\to$ melanesia the oracle now recovers $34\%$ of the
transfer gap, against the $72\%$ reported in an earlier draft from
pre-revision baselines. The change is attributable to corrected
baselines rather than to the oracle: the earlier source-only models
were trained with a contiguous-tail validation split, no early
stopping, and a normalisation defect that inflated the europe\_west
risk.

% =====================================================================
\section{Direct diagnosis of conditional shift}\label{sec:conditional}
% =====================================================================

The preliminary transfer results establish that adaptation fails, but
inference from method failure to shift type is indirect. We therefore
implement two direct diagnostics that measure whether $P(y\mid x)$
differs across domains without reference to any UDA method: a
model-free comparison of the conditional at matched inputs, and a
model-based comparison of the learnt mappings in the full input space.
The two are complementary---the first makes no modelling assumptions
but conditions only on a low-dimensional input projection, the second
uses the full input space but inherits the trained models' fitting
error. Agreement between them is the evidence.

\subsection{Matched-input conditional discrepancy}
\label{sec:matched_input}

To test Assumption~\ref{ass:causal} directly we compare the target
precipitation distribution between domains \emph{at matched inputs},
using observed samples only: no model, no training, no predictions.
Input space is partitioned into bins on the two dominant
precipitation-controlling variables, convective available potential
energy and total column water vapour, using shared quantile edges from
the pooled data so both domains occupy the same bins. Within each bin
$b$ --- that is, conditioning on approximately the same atmospheric
state --- we compute the exact empirical $W_1$ between the two domains'
log-precipitation distributions, and report the occupancy-weighted mean
$\overline{W}_1^{\text{cond}} = \sum_b w_b\, W_1(P_A(y\mid b),
P_B(y\mid b))$.

Under pure covariate shift (Assumption~\ref{ass:cov}) matched inputs
imply matched precipitation, so $\overline{W}_1^{\text{cond}}$ falls to
the level it would take were the two samples drawn from the same
distribution; under conditional shift the same atmospheric state
produces different precipitation and it remains well above that level.
We report $\overline{W}_1^{\text{cond}}$ itself as the statistic rather
than the ratio $r=\overline{W}_1^{\text{cond}}/W_1^{\text{uncond}}$
used in an earlier version of this diagnostic. That ratio has an
unstable denominator: two domains may have near-identical precipitation
marginals because a difference in $P(x)$ offsets a difference in
$P(y\mid x)$, in which case $W_1^{\text{uncond}}\to 0$ inflates $r$ at
perfectly adequate support. A large ratio is therefore not on its own
evidence of conditional shift, and we report it as context only.

\paragraph{Calibration requires a null.}
The empirical $W_1$ between two finite samples is positive even when
the distributions are identical, and positive again from the residual
covariate variation left inside each bin, since binning matches inputs
only to the bin width and only on the binned coordinates. Both inflate
$\overline{W}_1^{\text{cond}}$ by an unknown amount, so its raw value is
not interpretable on its own. We therefore split a single domain at
random into two halves and run the identical pipeline between them.
This yields the value the statistic takes under a true null, carrying
the same finite-sample noise, the same spatial correlation between
pixels, and the same within-bin residual covariate shift. We report
$\overline{W}_1^{\text{cond}}$ relative to that floor, and treat a
value not clearly exceeding it as no evidence of conditional shift. On
synthetic data with known ground truth the floor separates the regimes
cleanly: identical domains return $0.7\times$ the floor, constructed
covariate shift $2.1\times$, and constructed conditional shift
$137\times$.

\paragraph{Common support gates the comparison.}
The comparison is valid only where both domains populate the same bins.
We therefore report the common-support fraction, the share of each
domain's mass falling in bins both domains populate above a minimum
count. Low common support indicates the domains occupy disjoint input
regions, in which case $\overline{W}_1^{\text{cond}}$ is computed on an
unrepresentative overlap and is not interpretable as a conditional-shift
measurement.

Table~\ref{tab:matched_input} reports the diagnostic on the three
RainShift domains. \emph{All three pairs exhibit conditional shift},
differing in magnitude rather than in kind. The europe\_west to
horn-of-africa pair, at essentially complete common support, shows
$\overline{W}_1^{\text{cond}}$ at $87\times$ the null floor: these two
domains have the most similar input distributions of the three (the
smallest $W_1^{\text{uncond}}$) yet the largest conditional
discrepancy, direct evidence that the atmosphere-to-precipitation
mapping itself differs between the temperate and semi-arid-tropical
regimes. The horn-of-africa to melanesia pair, also at full support,
shows a smaller but unambiguous $14.5\times$ excess. Conditioning on
inputs does reduce its discrepancy substantially, from
$W_1^{\text{uncond}}=3.17$ to $\overline{W}_1^{\text{cond}}=0.40$;
what remains is nonetheless an order of magnitude above noise. We note
that under the ratio statistic this pair would have been classified as
predominantly covariate ($r=0.12$), because its unconditional
discrepancy is large, and that classification does not survive
calibration against the null. The europe\_west to melanesia pair is
flagged as low support ($44\%$): these domains occupy largely disjoint
regions of the (cape, tcw) plane, so no conditional statement can be
made --- a distinct barrier to covariate-shift adaptation, since input
marginals that do not overlap cannot be aligned.

\begin{table}[t]
\centering
\small
\begin{tabular}{lccccc}
\toprule
pair & $\overline{W}_1^{\text{cond}}$ & null floor & excess &
$W_1^{\text{uncond}}$ & support \\
\midrule
europe\_west $\leftrightarrow$ horn & $2.14$ & $0.025$ & $86.9\times$ &
$1.19$ & $99.6\%$ \\
horn $\leftrightarrow$ melanesia & $0.40$ & $0.027$ & $14.5\times$ &
$3.17$ & $100\%$ \\
europe\_west $\leftrightarrow$ melanesia & $2.15$ & $0.029$ &
($74.5\times$) & $1.98$ & $43.8\%$ \\
\bottomrule
\end{tabular}
\caption{Matched-input conditional discrepancy at $5\times5$ bins over
$600$ frames per domain. The statistic is
$\overline{W}_1^{\text{cond}}$; the null floor is the same statistic
computed between two random halves of a single domain, and excess is
their ratio. The europe\_west to melanesia row is support-limited and
its excess is not interpretable. Verdicts are unchanged under an
$8\times8$, $300$-frame configuration (excesses $37.6\times$,
$7.8\times$), and $\overline{W}_1^{\text{cond}}$ itself is stable
across the two ($2.141$ vs $2.140$; $0.497$ vs $0.395$): the coarser
binning changes the null, not the signal.}
\label{tab:matched_input}
\end{table}

\paragraph{Two limitations bound what this can claim.}
Because the binning conditions on a two-dimensional projection, the
errors are one-sided: a low value is conservative, since conditioning
on further variables can only reduce apparent conditional shift,
whereas a high value is an upper bound, since a difference in an
unbinned variable appears as different precipitation at nominally
matched inputs. Separately, the test is pixel-wise --- inputs and
target are coarsened to a common grid and pixels treated as
exchangeable --- so it measures $P(y_{\text{pixel}}\mid
\text{cape}_{\text{pixel}}, \text{tcw}_{\text{pixel}})$, whereas the
downscaling model learns $P(Y_{\text{field}}\mid X_{\text{field}})$, in
which precipitation at a pixel depends on the surrounding atmospheric
state. Two domains could share the field-wise conditional and still
differ pixel-wise. Pixels within a frame are also strongly correlated,
so the effective sample size per bin is well below the raw count; the
null baseline absorbs this, since it inherits the same correlation
structure. Both limitations motivate the model-based diagnostic below,
which operates on fields in the full nine-channel input space.

\subsection{Model-based conditional discrepancy}
\label{sec:model_based}

The source-only model $f_A$ trained on domain $A$ is an estimate of
$P_A(y\mid x)$. We use the set of trained source-only models to compare
conditionals in the full nine-channel input space, free of binning. For
an ordered pair $(A,B)$ we report two quantities. The transfer penalty
$\mathcal{R}_T(f_A\!\to\!B) - \mathcal{R}_B(f_B\!\to\!B)$ measures how
much worse $A$'s mapping is on $B$ than $B$'s own model, isolating the
mapping difference from $B$'s irreducible error. The model disagreement
is the standardised mean squared difference between $f_A$ and $f_B$
predictions on identical $B$ inputs; because both models receive the
same inputs, any disagreement is a pure difference in the learnt
mapping, with no ground-truth or target-side extrapolation term. A
concentration statistic, the ratio of the ninetieth percentile to the
median of the per-sample disagreement, distinguishes a diffuse mapping
difference (genuine conditional shift) from disagreement concentrated
in a tail of inputs (a support edge, i.e. extrapolation rather than
conditional shift).

Across all six ordered pairs the model disagreement is substantial
($0.22$ to $0.85$ in standardised units) and, by the concentration
statistic ($1.9$ to $3.6$, far below the tail threshold), diffuse
rather than support-concentrated. No pair is a support-edge artefact,
including europe\_west to melanesia, whose diffuse disagreement in the
full input space complements the low-support finding of the
matched-input test: even where the domains fail to overlap in the
(cape, tcw) projection, the extrapolated mappings differ throughout.
The disagreement magnitudes order by climatological distance---the two
tropical domains are most similar (melanesia to horn-of-africa,
$0.22$), europe\_west is most distinct from both ($0.65$ to
$0.85$)---mirroring the matched-input result, including the
europe\_west to horn-of-africa conditional reading, which is thereby
confirmed in the full input space rather than being an artefact of the
two-variable projection.

We caution that the transfer penalty is corrupted where the source-only
baseline is itself weak: for europe\_west as target the own-model risk
is anomalously high, producing negative penalties in which a foreign
model appears to outperform the home model. We therefore report the
model disagreement, which is robust to a weak baseline as it does not
reference the baseline's own risk, as the primary conditional-shift
statistic, and defer penalty-based quantities to the revised baselines.

\subsection{Synthesis}

Both diagnostics, applied independently and by different mechanisms,
agree that the RainShift domains differ in $P(y\mid x)$ and that the
magnitude of that difference orders by climatological distance. This is
a direct measurement of conditional shift rather than an inference from
adaptation failure, and it identifies which of the identifiability
assumptions holds per pair: Assumption~\ref{ass:cov} (covariate shift)
is tenable for \emph{no} pair in this triple. Calibrating the
matched-input statistic against a within-domain null removes the one
apparent exception --- horn-of-africa to melanesia appears
predominantly covariate under a ratio statistic, but retains a
$14.5\times$ excess over the noise floor once that floor is measured ---
while the europe\_west to melanesia pair is barred by non-overlapping
support rather than by a measured conditional difference. All three
therefore lie beyond the reach of marginal-alignment methods, regardless
of weighting. The result explains the observed adaptation failure at its
cause. It also withdraws a differential prediction that the
uncalibrated diagnostic appeared to license: with no pair classified as
covariate-dominated, we do not expect covariate alignment to recover
budget anywhere, and the joint-distribution method is the only candidate
in the taxonomy that remains consistent with the measured shift.

% =====================================================================
\section{Physical evaluation}\label{sec:physical}
% =====================================================================

Aggregate metrics describe how much a UDA method helps but not why or
where. RainShift's nine input channels capture distinct physical
processes (thermodynamic instability via CAPE, moisture availability
via total column water vapour, cloud mass via TCLW, synoptic dynamics
via wind components and surface pressure), which makes it possible to
stratify target performance by physical regime.

We bin target inputs into regimes defined by two- or three-variable
cuts of input space (e.g. CAPE--humidity, surface-pressure-anomaly--%
wind-magnitude) and report per-regime target error for the source-only
baseline and each UDA method. The expected output is a regime-by-method
matrix that reveals which physical processes each method helps or
hinders. Methods that improve performance uniformly across regimes are
acting as generic regularisers; methods that improve specifically in
convective regimes are addressing the tropical-precipitation shift
between europe\_west and melanesia; methods that worsen performance
in synoptic regimes are damaging the parts of the transfer that
source-only training already handled correctly.

The regime stratification is paired with the per-regime oracle
$\widehat{\lambda}_k$ estimate from Section~\ref{sec:lambda_k},
yielding a falsifiable test of the framework: if regime $k$ has
$\lambda_k$ large, no covariate-shift UDA method should improve target
performance specifically on regime $k$. Methods that violate this
prediction either fail Assumption~\ref{ass:inv} (in which case they
must be working by an orthogonal mechanism, e.g. regularisation that
helps in some regimes by accident) or signal that the regime
partition is itself misaligned with the true conditional shift
structure.

% =====================================================================
\section{Expected contributions}
% =====================================================================

The project will produce four contributions:

\textit{A pre-training diagnostic framework for UDA evaluation},
combining joint-marginal alignment residuals, multi-source disagreement
on supported inputs, and physical-prior consistency checks, each
mapped to a specific identifiability assumption
(Section~\ref{sec:identifiability}) and validated against the
post-hoc oracle measurement of $\lambda$. This framework is the
project's most generalisable methodological contribution; it applies
to any transfer-learning setting where practitioners must decide
whether to commit compute to UDA before observing target labels.

\textit{A novel regime-decomposed irreducible risk
$\lambda_k$} (Section~\ref{sec:lambda_k}) that localises conditional
shift to specific physical regimes of input space, with a
corresponding decomposition result (Proposition~\ref{prop:decomp})
and an empirical estimator via per-regime oracle joint training. To
our knowledge this is the first treatment of regime-decomposed
irreducible risk for regression with physically meaningful continuous
partitions, and constitutes the project's most original theoretical
contribution.

\textit{A systematic UDA benchmark on RainShift} across input, feature,
output, and inference-level methods for both deterministic and
probabilistic super-resolution architectures, evaluated on a
$3\times 3$ transfer matrix with multi-seed confidence intervals. The
benchmark explicitly reports source-MSE inflation alongside target
metrics to expose regression-to-mean artefacts in feature-moment
methods.

\textit{A physical-regime-stratified evaluation of UDA method success
and failure}, identifying which atmospheric processes each method
addresses and which it leaves untouched, with the regime partition
constructed from RainShift input channels alone and therefore
applicable in target-label-free deployment.

% =====================================================================
\section{Future directions}\label{sec:future}
% =====================================================================

Three theoretical extensions are natural follow-ups beyond the scope
of the present project.

\emph{PAC-Bayesian domain adaptation bounds.} The bounds in
Section~\ref{sec:theory} are worst-case over the hypothesis class.
PAC-Bayes-DA bounds~\citep{germain2013pac} replace the worst-case
complexity term with one that depends on a posterior over weights,
which for deep networks can be empirically
non-vacuous~\citep{dziugaite2017computing}. AFM's flow network
provides a natural posterior approximation via stochastic predictions,
which could be exploited to compute non-vacuous bounds on
$\widehat\lambda$ for the case studies.

\emph{Causal-mechanism transferability.} Assumption~\ref{ass:causal}
is a strong claim about which inputs constitute a sufficient causal
set. A more refined treatment would identify, per regime, which
sub-grid processes are the dominant causal omissions, with the
regime-decomposed $\lambda_k$ providing the empirical signal. This
connects naturally to recent work on causal-mechanism-aware
generalisation~\citep{scholkopf2012causal}.

\emph{Tighter joint-distribution shift estimators.} The Joint-OT
quantity $W_1(\mathcal{D}_S, \mathcal{D}_T)$ provides a single scalar
that captures both covariate and conditional shift, but is not
empirically estimable without target labels. Surrogate estimators of
the joint shift, perhaps via the regime-decomposed $\lambda_k$ as an
intermediary, could close this gap.

% =====================================================================
\bibliographystyle{unsrtnat}
\bibliography{references}

\end{document}